\documentclass[a4paper]{report}
\usepackage[utf8]{inputenc}
\usepackage[portuguese]{babel}
\usepackage{hyperref}
\usepackage{a4wide}
\hypersetup{pdftitle={PCP - OpenMP},
pdfauthor={João Teixeira, José Ferreira},
colorlinks=true,
urlcolor=blue,
linkcolor=black}
\usepackage{subcaption}
\usepackage{listings}
\usepackage{booktabs}
\usepackage{multirow}
\usepackage{appendix}
\usepackage{tikz}
\usepackage{authblk}
\usepackage{bashful}
\usepackage{verbatim}
\usepackage{amssymb}
\usepackage{multirow}
\usepackage{mwe}
\usepackage[parfill]{parskip}
\usetikzlibrary{positioning,automata,decorations.markings}
\AfterEndEnvironment{figure}{\noindent\ignorespaces}
\AfterEndEnvironment{table}{\noindent\ignorespaces}

\begin{document}

\title{Paradigmas de Computação Paralela\\Bucket Sort com OpenMP}
\author{João Teixeira (A85504) \and José Filipe Ferreira (A83683)}
\date{\today}

\begin{center}
    \begin{minipage}{0.75\linewidth}
        \centering
        \includegraphics[width=0.4\textwidth]{images/eng.jpeg}\par\vspace{1cm}
        \vspace{1.5cm}
        \href{https://www.uminho.pt/PT}
        {\color{black}{\scshape\LARGE Universidade do Minho}} \par
        \vspace{1cm}
        \href{https://www.di.uminho.pt/}
        {\color{black}{\scshape\Large Departamento de Informática}} \par
        \vspace{1.5cm}
        \maketitle
    \end{minipage}
\end{center}

\tableofcontents

\pagebreak

\chapter{Introdução}
O algoritmo escolhido para o projeto da unidade curricular de Computação
Paralela e Distribuída foi o \textit{Bucket Sort}.

Inicialmente desenvolvemos uma versão sequencial do projeto e procedemos ao
\textit{benchmarking} do programa resultante. Em seguida convertemos a
implementação sequencial numa versão com utilização de memoria partilhada
fazendo uso de \textit{OpenMP} e comparamos o resultado com a versão sequencial
desenvolvida anteriormente.

Ao longo deste relatório iremos descrever a metodologia utilizada e os
resultados de \textit{benchmarking} obtidos ao longo deste projeto.

De notar que todos os benchmarks descritos foram efeutados em nós do tipo 641 do
cluster \textit{SeARCH}, e todos os executáveis foram compilados o
\textit{MPIcc} na versão 1.8.1 do \textit{OpenMPI}, utilizando a versão
7.2.0 do \textit{gcc}, e com as flags \textit{-O3 -std=c11} para garantir 
a maior consistência entre os diferentes benchmarks. Na
ausência de indicação, os benchmarks foram efetuados com um input de 100000000
de elementos aleatórios, entre -1000 e 500000.

\chapter{OpenMPI} \label{chap:ompi}

O algoritmo de Bucket Sort adaptado para \textit{MPI} divide-se em quatro fases,
divisão dos elementos em chunks para serem enviados para os diferentes
processos, separação dos elementos pelos diversos baldes, junção e ordenação dos
baldes dos diferentes processos, e por fim a agregação dos baldes ordenados num
array final.

A nossa implementação realiza a primeira fase de divisão dos elementos em chunks
num unico processo, onde é difundido o máximo e o minímo dos elementos a
ordenar, e o tamanho dos buffers que os diferentes processos devem ter para
garantir a boa receção do chunk que lhe corresponde. Depois desta difusão, é
enviado a todos os processos o chunk correspondendo, realizando P envios, um por
cada processo, com o tamanho máximo de (P/N + 1) * 4 bytes, onde P é o número de
processos e o N corresponde ao número de elementos total a ordenar.

Recebido pelos processos os chunks respetivos, cada um vai percorrer os
elementos recebidos, atribuindo-o cada um ao respetivo balde. Depois de todos os
baldes preenchidos, cada processo vai enviar cada balde para um processo que
irá juntar ao mesmo balde vindo dos outros processos. Esta distribuição é
efetuada para que todos os processos recebam um número equilibrado de
buckets, onde cada processo irá receber B/P baldes, caso existam mais baldes
que processos, senão irá receber ou um ou zero baldes, dependendo do seu id. 
Aqui cada processo irá realizar B envios, com um máximo (P/N + 1) * 4 bytes, 
com B sendo o número de baldes, e o P e N indicados acima.

Finda a receção de todos os buckets que um processo tem para processar, estes
irão proceder à ordenação destes, e finda esta são agrupados num só array os
baldes consecutivos e ordenados. Depois deste passo, é enviado este array para o
processo que tinha inicialmente feito a distribuição dos chunks. É realizado um
envio no máximo com N * 4 bytes.

Neste processo que recebeu os arrays de todos os processos que trataram de algum
bucket, são agrupados e retornado o resultado final.

\appendix

\chapter{Benchmarking}
\begin{table}[h]
    \centering
    \begin{tabular}{|l|l|l|l|l|}
        \hline
        \multicolumn{5}{|c|}{Times}                                      \\ \hline
        Threads & Old Critical & Old Task    & New       & New Collapsed \\ \hline
        Seq     & 16,209999    & 16,209999   & 22,299999 & 22,299999     \\ \hline
        2       & 43,23        & 305,169983  & 11,403177 & 11,968859     \\ \hline
        4       & 100,150002   & 937,929993  & 6,880816  & 6,872915      \\ \hline
        8       & 187,929993   & 2227,460205 & 4,622566  & 5,965684      \\ \hline
        16      & 479,750031   & 5547,610352 & 2,362948  & 4,678829      \\ \hline
        32      & 1213,959961  & 12624,78906 & 2,363651  & 4,859851      \\ \hline
        64      & 1726,070068  & 11589,56934 & 2,449052  & 4,059251      \\ \hline
    \end{tabular}
    \caption{\label{tab:Times}Tempos de execução iniciais, com 10 baldes}
\end{table}

\begin{table}[h]
    \centering
    \begin{tabular}{|l|l|l|l|l|l|}
        \hline
        Threads & 5 buckets & 10 buckets & 20 buckets & 40 buckets & 80 bucket \\ \hline
        2       & 12.520761 & 11.218051  & 12.921473  & 17.016363  & 25.398283 \\ \hline
        4       & 8.365767  & 6.739361   & 6.511606   & 8.512317   & 12.743460 \\ \hline
        8       & 5.218658  & 4.984042   & 3.952061   & 4.323404   & 6.438317  \\ \hline
        16      & 5.344770  & 3.167999   & 3.042659   & 2.632886   & 3.465635  \\ \hline
        32      & 5.248249  & 3.147622   & 2.148584   & 2.392601   & 2.836905  \\ \hline
        64      & 5.270648  & 3.161402   & 2.366316   & 2.002006   & 3.260201  \\ \hline
    \end{tabular}
    \caption{\label{tab:fix_size}Tempos de execução, fazendo variar o
    número de baldes utilizados}
\end{table}

\begin{table}[h!]
    \centering
    \begin{tabular}{|l|l|l|l|l|l|}
        \hline
        Threads & n\_threads & 2*n\_threads & 4*n\_threads & 8*n\_threads & 16*n\_threads \\ \hline
        2       & 11.395469  & 10.310162    & 10.808166    & 12.206050    & 15.356203     \\ \hline
        4       & 6.129920   & 5.490817     & 6.108120     & 7.681805     & 11.012663     \\ \hline
        8       & 3.719489   & 3.100843     & 3.883367     & 5.560686     & 9.027642      \\ \hline
        16      & 2.217584   & 2.001313     & 2.854720     & 4.599587     & 8.157944      \\ \hline
        32      & 1.824530   & 2.292776     & 3.511230     & 6.276313     & 11.743461     \\ \hline
        64      & 2.486487   & 3.223805     & 5.932982     & 10.554493    & 20.222078     \\ \hline
    \end{tabular}
    \caption{\label{tab:var_size}Tempos de execução, fazendo variar o
    número de baldes, com base no número de threads.}
\end{table}

\begin{table}[h]
    \centering
    \begin{tabular}{|l|l|l|l|l|l|l|}
        \hline
        Threads & 30       & 32       & 34       & 36       & 38       & 40       \\ \hline
        16      & 2.085598 & 2.008397 &          &          &          &          \\ \hline
        32      & 1.846964 & 1.813501 & 2.572607 & 2.442342 & 2.408187 & 2.349852 \\ \hline
        40      & 1.907169 & 1.799508 & 1.779662 & 1.718217 & 1.788759 & 1.881975 \\ \hline
        64      & 2.599405 & 2.536593 & 2.106488 & 1.911450 & 1.976940 & 1.977957 \\ \hline
    \end{tabular}
    \caption{\label{tab:cutoff}Tempos de execução, com o número de baldes a ser o
    dobro do número de threads, e variando o número máximo de baldes
    utilizados.}
\end{table}

\begin{table}[h]
    \centering
    \begin{tabular}{|l|l|l|}
        \hline
        \multicolumn{1}{|c|}{Threads} & \multicolumn{1}{c|}{Time} & \multicolumn{1}{c|}{Speedup} \\ \hline
        Seq                           & 20,629999                 & 1                            \\ \hline
        2                             & 10,346263                 & 1,993956562                  \\ \hline
        4                             & 5,57325                   & 3,701610192                  \\ \hline
        8                             & 3,288632                  & 6,273124813                  \\ \hline
        16                            & 2,199246                  & 9,380487222                  \\ \hline
        32                            & 1,677771                  & 12,29607557                  \\ \hline
        64                            & 1,814721                  & 11,36813813                  \\ \hline
    \end{tabular}
    \caption{\label{tab:final}Tempos finais de execução e respetivo speedup, com a
    versão sequencial a utilizar 4 baldes.}
\end{table}

\begin{table}[h]
    \centering
    \begin{tabular}{|l|l|l|}
        \hline
        \multicolumn{1}{|c|}{Threads} & \multicolumn{1}{c|}{Time} & \multicolumn{1}{c|}{Speedup} \\ \hline
        Seq                           & 1,57                      & 1                            \\ \hline
        2                             & 1,485135                  & 1,057142953                  \\ \hline
        4                             & 1,49086                   & 1,053083455                  \\ \hline
        8                             & 1,498286                  & 1,047864026                  \\ \hline
        16                            & 1,504521                  & 1,043521493                  \\ \hline
        32                            & 1,45915                   & 1,075968886                  \\ \hline
        64                            & 1,456144                  & 1,078190069                  \\ \hline
    \end{tabular}
    \caption{\label{tab:unbal}Tempos de execução para 10000000 elementos aleatórios entre -100 e 100,
    com um elemento sendo 500000}
\end{table}
\end{document}
